\documentclass[11pt,a4paper]{article}
\usepackage{shared/luxcover}
\usepackage[utf8]{inputenc}
\usepackage[T1]{fontenc}
\usepackage{amsmath,amssymb,amsfonts,amsthm}
\usepackage{graphicx}
\usepackage{hyperref}
\usepackage{booktabs}
\usepackage{algorithm}
\usepackage{algpseudocode}
\usepackage{enumitem}
\usepackage{xcolor}

\hypersetup{colorlinks=true,linkcolor=black,citecolor=blue,urlcolor=blue}

\newtheorem{theorem}{Theorem}[section]
\newtheorem{lemma}[theorem]{Lemma}
\newtheorem{definition}{Definition}[section]
\newtheorem{remark}{Remark}[section]

\title{LP-308: QuantumSwap\\[4pt]
\large A Post-Quantum AMM Protocol with Mixed\\
Fungible/Non-Fungible Liquidity Pools}

\author{
Zach Kelling\thanks{Corresponding author: zach@lux.network}\\
\textit{Lux Industries, Inc.}\\
\texttt{research@lux.network}
}

\date{Original: March 2024\quad|\quad Revised: April 2026}

\begin{document}

\luxcoverpage

\begin{abstract}
QuantumSwap is an automated market maker protocol designed for the
Lux Network with three distinguishing properties: post-quantum
encryption of order flow via TFHE, mixed fungible/non-fungible token
liquidity pools with configurable bonding curves, and a multi-chain
teleport protocol using zero-knowledge bridges with Lamport and
lattice-based signatures.

This paper formalizes the bonding curve algorithms (linear,
exponential, XYK), the interest rate mechanism for RWA-backed
price-stable currencies, and the capital efficiency model. We provide
the concrete interest rate schedule and analyze its stability
properties.

Section~\ref{sec:current} maps the QuantumSwap design to the current
Lux exchange architecture: \texttt{luxfi/dex} (CLOB engine),
\texttt{luxfi/standard/contracts/amm/} (AMM V1--V4 contracts), and
the USDL stablecoin.
\end{abstract}

\tableofcontents
\newpage

%% ============================================================
\section{Introduction}

\subsection{Motivation}

Automated market makers revolutionized decentralized exchange by
replacing order books with deterministic pricing functions. However,
existing AMMs suffer from three limitations:

\begin{enumerate}[nosep]
  \item \textbf{No privacy}: All orders are visible on-chain before
        execution, enabling front-running and sandwich attacks.
  \item \textbf{Fungible tokens only}: NFTs and real-world asset tokens
        cannot participate in standard AMM pools.
  \item \textbf{Single-chain}: Cross-chain liquidity requires third-party
        bridges with variable security properties.
\end{enumerate}

QuantumSwap addresses all three: homomorphic encryption prevents order
visibility, mixed FT/NFT pools extend liquidity to non-fungible assets,
and a native teleport protocol provides cross-chain token movement with
post-quantum security.

\subsection{Historical Context}

\begin{center}
\begin{tabular}{lll}
\toprule
\textbf{Year} & \textbf{Protocol} & \textbf{Innovation} \\
\midrule
2013 & BitShares & Decentralized order books \\
2016 & 0x Protocol & Gas-efficient meta-transactions \\
2018 & Uniswap V1 & First AMM, LP tokens, ETH/ERC-20 pairs \\
2020 & Uniswap V2 & ERC-20/ERC-20 pairs, flash swaps, oracle feeds \\
2021 & Uniswap V3 & Concentrated liquidity, range orders~\cite{uniswapv3} \\
2024 & QuantumSwap & Mixed FT/NFT pools, FHE-encrypted order flow \\
\bottomrule
\end{tabular}
\end{center}

%% ============================================================
\section{Bonding Curve Algorithms}

QuantumSwap pools are parameterized by a bonding curve that determines
price dynamics. Each \texttt{LUXQSAPair} contract is associated with
exactly one curve. The curves are pure functions: they do not modify
pool state. State updates occur in the \texttt{LUXRWAPair} contract
after the curve is consulted.

\subsection{Notation}

Let $\mathit{spot}$ denote the current spot price (instantaneous price
to sell one item to the pool), $\delta$ the curve parameter, and $x$
the trade size. The spot price to \emph{buy} one item from the pool is
always $\mathit{spot} + \delta$ (additive) or $\mathit{spot} \cdot
(1 + \delta)$ (multiplicative), depending on the curve.

\subsection{Linear Bonding Curve}

The linear curve adjusts price by a constant additive amount $\delta$
per item traded:

\begin{align}
\mathit{spot}_{\text{after sell}} &= \mathit{spot} - \delta
  \label{eq:linear-sell} \\
\mathit{spot}_{\text{after buy}} &= \mathit{spot} + \delta
  \label{eq:linear-buy}
\end{align}

For a multi-item sell of $x$ items, the total payout is:

\begin{equation}\label{eq:linear-multi}
\text{Payout}(x) = \sum_{i=0}^{x-1} (\mathit{spot} - i\delta)
= x \cdot \mathit{spot} - \delta \cdot \frac{x(x-1)}{2}
\end{equation}

\begin{remark}
The LP must set $\delta$ with the same decimal precision as the
pair's underlying token. For an ETH-denominated pair with
$\mathit{spot} = 1\,\text{ETH}$ and $\delta = 0.1\,\text{ETH}$,
selling 5~NFTs yields
$5 \cdot 1 - 0.1 \cdot 10 = 4\,\text{ETH}$, distributed as
$\{1, 0.9, 0.8, 0.7, 0.6\}$.
\end{remark}

\subsection{Exponential Bonding Curve}

The exponential curve adjusts price by a multiplicative factor:

\begin{align}
\mathit{spot}_{\text{after buy}} &= \mathit{spot} \cdot (1 + \delta)
  \label{eq:exp-buy} \\
\mathit{spot}_{\text{after sell}} &= \frac{\mathit{spot}}{1 + \delta}
  \label{eq:exp-sell}
\end{align}

After $x$ consecutive buys from an initial spot price $p_0$:

\begin{equation}\label{eq:exp-multi}
\mathit{spot}(x) = p_0 \cdot (1+\delta)^x
\end{equation}

The total cost of buying $x$ items is:

\begin{equation}
\text{Cost}(x) = p_0 \cdot (1+\delta) \cdot
\frac{(1+\delta)^x - 1}{\delta}
\end{equation}

\begin{remark}
With $p_0 = 2\,\text{ETH}$ and $\delta = 0.5$ (50\%), the price
sequence for 3~buys is $\{3, 4.5, 6.75\}$ ETH. Selling into the
pool divides by $1.5$ at each step.
\end{remark}

\subsection{XYK Bonding Curve}

The XYK curve maintains a constant-product invariant over virtual
reserves, adapted for NFT pools. At pool creation:

\begin{align}
R_{\text{nft}} &= \text{(number of items to buy or sell)} + 1
  \label{eq:xyk-nft} \\
R_{\text{token}} &= R_{\text{nft}} \cdot \mathit{spot}
  \label{eq:xyk-token}
\end{align}

The price to buy $x$ items from the pool (receiving NFTs, sending
tokens) is:

\begin{equation}\label{eq:xyk-buy}
\text{Cost}(x) = \frac{x \cdot R_{\text{token}}}{R_{\text{nft}} - x}
\end{equation}

The price to sell $x$ items to the pool is:

\begin{equation}\label{eq:xyk-sell}
\text{Payout}(x) = \frac{x \cdot R_{\text{token}}}{R_{\text{nft}} + x}
\end{equation}

After each trade, the virtual reserves update:
\begin{itemize}[nosep]
  \item Pool sells NFTs:
    $R_{\text{nft}} \leftarrow R_{\text{nft}} - x$, \;
    $R_{\text{token}} \leftarrow R_{\text{token}} + \text{Cost}(x)$
  \item Pool buys NFTs:
    $R_{\text{nft}} \leftarrow R_{\text{nft}} + x$, \;
    $R_{\text{token}} \leftarrow R_{\text{token}} - \text{Payout}(x)$
\end{itemize}

A concentration parameter allows LPs to tighten or loosen the curve
around the current price, analogous to Uniswap V3's concentrated
liquidity~\cite{uniswapv3} but applied to discrete (NFT) assets.

%% ============================================================
\section{Mixed FT/NFT Liquidity Pools}

\subsection{Pool Types}

QuantumSwap supports three pool configurations:

\begin{center}
\begin{tabular}{lll}
\toprule
\textbf{Type} & \textbf{Assets} & \textbf{Use Case} \\
\midrule
FT/FT & Two ERC-20 tokens & Standard AMM \\
FT/NFT & ERC-20 + ERC-721 & RWA token trading \\
NFT/NFT & Two ERC-721 collections & Cross-collection swaps \\
\bottomrule
\end{tabular}
\end{center}

\subsection{FT/NFT Mechanics}

In an FT/NFT pool, the NFT side is treated as a discrete quantity
while the FT side is continuous. The bonding curve operates on the
NFT count, producing FT prices. This enables trade sequences like:

\[
\text{FT} \xrightarrow{\text{Pool A}} \text{NFT}_1
\xrightarrow{\text{Pool B}} \text{NFT}_2
\xrightarrow{\text{Pool C}} \text{FT}'
\]

Each hop uses the bonding curve of the respective pool. The router
contract finds the optimal path across available pools.

\subsection{Capital Efficiency}

Concentrated liquidity in QuantumSwap achieves up to $4000\times$
capital efficiency over Uniswap V2's uniform distribution, consistent
with the analysis in Adams et al.~\cite{uniswapv3}. The improvement
comes from restricting liquidity to a narrow price range rather than
distributing it uniformly from zero to infinity.

\begin{remark}
The $4000\times$ figure is a theoretical maximum for a single tick
range. Practical improvement depends on the LP's price range selection
relative to realized volatility. Empirical data from Uniswap V3
shows typical improvements of $4\times$ to $8\times$ for actively
managed positions.
\end{remark}

%% ============================================================
\section{RWA-Backed Price-Stable Currencies}

\subsection{Design}

A Real World Asset (RWA) token represents a claim on a physical
asset (gold, silver, uranium, real estate). When users stake RWA tokens
into a governance pool, they mint a price-stable currency pegged to
the underlying asset's market value. The bonding curve automatically
adjusts an interest rate to maintain the peg.

\subsection{Interest Rate Algorithm}\label{sec:interest}

Let:
\begin{itemize}[nosep]
  \item $b_{\text{LUX}}$ = value of one unit of the RWA token on
        QuantumSwap (in USD)
  \item $b_{\text{market}}$ = value of the same underlying asset on
        major external exchanges (in USD)
  \item $\text{premium} = (b_{\text{LUX}} - b_{\text{market}}) /
        b_{\text{market}}$
\end{itemize}

The annualized interest rate $r$ is determined by:

\begin{equation}\label{eq:interest-rate}
r = \begin{cases}
  0\%   & \text{if } \text{premium} = 0\% \\
  2\%   & \text{if } 0\% < \text{premium} \leq 2\% \\
  \text{premium} & \text{if } 2\% < \text{premium} \leq 15\% \\
  15\%  & \text{if } \text{premium} > 15\%
\end{cases}
\end{equation}

\subsection{Stability Analysis}

The rate schedule in~\eqref{eq:interest-rate} has three regimes:

\begin{enumerate}[nosep]
  \item \textbf{Dead zone} ($\text{premium} = 0$): No incentive to
        issue or redeem. The peg is exact.
  \item \textbf{Threshold} ($0 < \text{premium} \leq 2\%$): A fixed
        2\% rate provides a clear signal to issuers. The step function
        avoids oscillation around zero.
  \item \textbf{Linear tracking} ($2\% < \text{premium} \leq 15\%$):
        Rate tracks premium 1:1, providing proportional incentive.
  \item \textbf{Cap} ($\text{premium} > 15\%$): Protects holders from
        unbounded interest obligations.
\end{enumerate}

\begin{remark}
The 2\% threshold was chosen empirically: users respond more
reliably to a discrete step than to continuous micro-incentives. The
15\% cap prevents death spirals where rising rates cause holder panic
and further depegging.
\end{remark}

\subsection{Comparison to Prior Art}

BitShares' pegged currencies (bitUSD, bitCNY) were collateralized
solely by BTS, a volatile crypto-asset. This created reflexive
collapse risk: BTS price drops $\rightarrow$ collateral falls
$\rightarrow$ peg breaks $\rightarrow$ BTS price drops further.
QuantumSwap's design requires 1:1 physical asset collateral (gold,
uranium, etc.), breaking the reflexivity. The interest rate mechanism
then handles \emph{demand}-side deviations from peg, while the
physical collateral handles \emph{supply}-side solvency.

%% ============================================================
\section{Privacy: Homomorphic Encryption of Order Flow}

\subsection{Threat Model}

In a standard AMM, pending transactions are visible in the mempool.
An adversary can:
\begin{itemize}[nosep]
  \item \textbf{Front-run}: Place a transaction before the victim's to
        profit from the price impact.
  \item \textbf{Sandwich}: Place transactions before and after the
        victim's, extracting value from the spread.
\end{itemize}

\subsection{TFHE-Protected Pools}

QuantumSwap encrypts order parameters (amount, direction, slippage
tolerance) under TFHE before submission. The AMM contract evaluates
the bonding curve on encrypted inputs:

\begin{itemize}[nosep]
  \item Comparison operations (is the price within slippage tolerance?)
        execute as TFHE boolean circuits.
  \item Reserve updates use TFHE integer arithmetic.
  \item The result (trade executed or rejected) is encrypted; only
        the trader can decrypt their outcome.
\end{itemize}

The performance cost is significant: a single TFHE-protected swap
requires approximately 100 bootstrapping operations ($\sim$10 seconds
at PN10QP27 parameters). This is acceptable for high-value RWA
trades but not for high-frequency token swaps.

%% ============================================================
\section{Multi-Chain Teleport Protocol}

\subsection{Architecture}

The teleport protocol enables cross-chain token movement without
relying on third-party bridges. It consists of:

\begin{enumerate}[nosep]
  \item \textbf{Source chain lock}: Tokens are locked in a vault
        contract on the source chain.
  \item \textbf{ZK proof generation}: A zero-knowledge proof of the
        lock event is generated, attesting to the token identity,
        amount, and destination without revealing the sender.
  \item \textbf{Relay}: A trustless relayer transmits the proof across
        chains. The relayer cannot modify the proof or redirect the
        tokens.
  \item \textbf{Destination chain mint}: The proof is verified
        on-chain using a ZK verifier (Groth16 or PLONK). On
        verification, equivalent tokens are minted on the destination.
\end{enumerate}

\subsection{Post-Quantum Security}

The ZK bridge uses Lamport signatures for proof authentication (quantum
resistance from hash-only construction) and lattice-based encryption
for the relay channel. The combination provides post-quantum security
for the bridge even if the underlying chains use ECDSA.

%% ============================================================
\section{Gas-Free Meta-Transactions}

QuantumSwap implements ERC-2771~\cite{erc2771} compliant
meta-transactions. Users sign transactions off-chain; a trusted
forwarder contract verifies the signature and forwards the call to the
AMM, paying gas on the user's behalf. This eliminates the gas barrier
for new users.

Formally, a meta-transaction $M = (\mathit{to}, \mathit{value},
\mathit{data}, \mathit{sig})$ is a standard transaction with the sender
address derived from $\mathit{sig}$ rather than from \texttt{msg.sender}.
The forwarder verifies $\mathit{sig}$ against a nonce to prevent replay.

%% ============================================================
\section{Current Status}\label{sec:current}

\subsection{Architecture Evolution}

The Lux exchange architecture has evolved from AMM-first
(QuantumSwap, 2024) to a hybrid CLOB + AMM design:

\begin{center}
\begin{tabular}{lll}
\toprule
\textbf{Component} & \textbf{2024 (QuantumSwap)} & \textbf{2026 (Production)} \\
\midrule
Matching engine & AMM bonding curves & CLOB (\texttt{luxfi/dex}) \\
AMM contracts & Proposed & Deployed (V1--V4) \\
Order privacy & TFHE on all orders & TFHE on Z-Chain only \\
Stablecoin & Gold/uranium-backed & USD-backed (USDL) \\
Cross-chain & ZK teleport & Warp messaging + MPC bridge \\
\bottomrule
\end{tabular}
\end{center}

\subsection{What Shipped}

\begin{itemize}[nosep]
  \item \textbf{AMM contracts} in \texttt{luxfi/standard/contracts/amm/}
        implement V1 (constant-product), V2 (multi-asset), V3
        (concentrated liquidity), and V4 (hook-extensible) designs.
        The bonding curve algorithms from Sections~2--3 are implemented
        in the V1/V2 pair contracts.
  \item \textbf{CLOB engine} (\texttt{luxfi/dex}) provides 434M ops/sec
        GPU-accelerated matching, used as the primary exchange engine.
        AMM pools serve as passive liquidity that the CLOB can route
        through.
  \item \textbf{USDL}: The stablecoin shipped as a straightforward
        USD-backed ERC-20, not the RWA-collateralized multi-asset
        design. Physical commodity backing (gold, uranium) was deferred
        due to regulatory complexity around commodity-backed tokens.
  \item \textbf{Mixed FT/NFT pools}: The \texttt{LUXRWAPair} contract
        is deployed for RWA token trading, including SecurityToken
        (ERC-20) swaps.
\end{itemize}

\subsection{What Changed}

\begin{itemize}[nosep]
  \item \textbf{CLOB over AMM}: For securities trading, a CLOB provides
        price discovery, order types (limit, stop, market), and
        regulatory compliance (best execution obligations) that AMMs
        cannot. AMM pools remain available for DeFi use cases.
  \item \textbf{Warp + MPC over ZK teleport}: The production bridge
        uses Warp messaging with BLS aggregation and 2-of-3 MPC
        threshold signing, which is simpler to audit and maintain than
        a full ZK bridge circuit.
  \item \textbf{USD-backed over commodity-backed}: USDL is 1:1
        USD-backed via ACH deposits. The interest rate mechanism
        from Section~\ref{sec:interest} was not needed for a
        fiat-backed stablecoin.
\end{itemize}

%% ============================================================
\section{Conclusion}

QuantumSwap introduced three ideas that influenced the production Lux
exchange: mixed FT/NFT liquidity pools with formal bonding curves,
FHE-encrypted order flow for front-running prevention, and an interest
rate mechanism for maintaining asset-backed pegs. The CLOB engine
superseded AMM as the primary matching mechanism for securities, while
AMM contracts remain deployed for DeFi liquidity provision. The
interest rate algorithm remains available for future commodity-backed
stablecoin deployments.

\bibliographystyle{plain}
\begin{thebibliography}{99}

\bibitem{uniswapv3}
H.~Adams, N.~Zinsmeister, M.~Salem, R.~Keefer, and D.~Robinson,
``Uniswap v3 Core,'' Uniswap Labs, 2021.

\bibitem{erc2771}
G.~S.~Vogelsteller et al., ``ERC-2771: Secure Protocol for Native
Meta Transactions,'' Ethereum Improvement Proposals, 2020.

\bibitem{sudoswap}
0xmons, ``sudoswap: An on-chain NFT AMM,''
\url{https://docs.sudoswap.xyz/}, 2022.

\bibitem{gentry2009}
C.~Gentry, ``A fully homomorphic encryption scheme,'' Ph.D.
dissertation, Stanford University, 2009.

\bibitem{lamport1979}
L.~Lamport, ``Constructing digital signatures from a one-way
function,'' Technical Report SRI-CSL-98, 1979.

\bibitem{lp105}
Z.~Kelling, ``LP-105: Quasar---Hybrid BLS + Lattice Consensus for
Post-Quantum Finality,'' Lux Industries, 2026.

\end{thebibliography}

\end{document}
